% Fragmento aditivo: no modifica el manuscrito fuente.
% Procedencia: U014, regla modal, incidencia y caracter multiplicativo.
% Este fragmento desarrolla el envolvente modal, no lo identifica con todo el TPK.

\subsection{Multiplicidad estructural del carácter de autoescala}

La publicación escalar identifica rutas distintas. En el envolvente modal
inducido por el calendario trítico puede determinarse exactamente esa
multiplicidad y conservarse la ruta mediante una representación reversible.

El bloque operativo $(+1,-1,0)$ publica la sucesión modal
$(\Sigma,\Pi,\Pi)$. Al contar sus transiciones se obtiene
\[
 F=\begin{pmatrix}0&1\\1&1\end{pmatrix},
 \qquad 0=\Sigma,\quad 1=\Pi.
\]
La ecuación del rayo positivo, $F(1,x)^{\mathsf T}=x(1,x)^{\mathsf T}$,
impone $x^2=x+1$. Su raíz positiva se reconoce como $\varphi$.
Esta construcción precede a la lectura escalar.

Se considera el conjunto de paseos dirigidos finitos
$\gamma=(x_0,\ldots,x_n)$ de esta incidencia. Se distinguen sus extremos
y longitudes, se admiten repeticiones y se incluyen los paseos triviales.
Este dominio es el envolvente modal; el estado completo del TPK conserva,
además, los registros que no aparecen en esta publicación.

El carácter multiplicativo documentado adopta la forma
\[
 \chi(\gamma)=\varphi^n\frac{r_{x_0}}{r_{x_n}}
 =\varphi^{k(\gamma)},\qquad
 r=(1,\varphi),\quad k(\gamma)=n+x_0-x_n.
\]
Las transiciones $\Sigma\to\Pi$, $\Pi\to\Pi$ y $\Pi\to\Sigma$ incrementan
$k$ en $0$, $1$ y $2$, respectivamente. El primer paso modifica la ruta
sin alterar su carácter. En particular, $(1,1)$ y $(0,1,1)$ publican
ambos $\varphi$.

\paragraph{Teorema de multiplicidad de publicación.}
Si $\mathcal F_k=\{\gamma:k(\gamma)=k\}$ y $N_k=|\mathcal F_k|$, entonces
\[
 N_0=3,\qquad
 N_k=2\operatorname{tr}(F^k)
 =2\bigl(\varphi^k+(-\varphi^{-1})^k\bigr)\quad(k\ge1).
\]
Por consiguiente,
\[
 \lim_{k\to\infty}N_k^{1/k}=\varphi.
\]

\paragraph{Demostración.}
Dados los extremos $i,j$, la longitud es $n=k-i+j$. Para $k\ge1$,
\[
 N_k=\operatorname{tr}(F^k)+(F^{k+1})_{01}+(F^{k-1})_{10}.
\]
La identidad $F^2=F+I$ da, para $n\ge1$,
\[
 F^n=\begin{pmatrix}f_{n-1}&f_n\\f_n&f_{n+1}\end{pmatrix},
 \qquad f_0=0,\quad f_1=1,\quad f_{n+1}=f_n+f_{n-1}.
\]
Los dos términos no diagonales anteriores suman
$\operatorname{tr}(F^k)$; el caso $k=1$ se obtiene con $F^0=I$.
Para $k=0$ sólo existen $(0)$, $(1)$ y $(0,1)$.
Los autovalores de $F$ son $\varphi$ y $-\varphi^{-1}$, lo que proporciona
la fórmula y su límite.

La autoescala admite así una caracterización adicional: es la tasa
exponencial de crecimiento de las genealogías modales finitas que
el carácter escalar hace indistinguibles.

\paragraph{Refinamiento por composición de los retornos.}
Sean $b(\gamma)$ el número de aristas $\Pi\to\Sigma$ y $c(\gamma)$ el de
aristas $\Pi\to\Pi$. Se cumple $k=2b+c$. La incidencia graduada
\[
 M(z,t)=\begin{pmatrix}0&1\\tz^2&z\end{pmatrix}
\]
registra simultáneamente exponente y retornos. No contiene ciclos de
grado cero. Por ello cada coeficiente cuenta un conjunto finito de paseos:
\[
 \sum_{k,b\ge0}N_{k,b}z^kt^b
 =\mathbf1^{\mathsf T}(I-M(z,t))^{-1}\mathbf1
 =\frac{3-z+tz^2}{1-z-tz^2}.
\]
Para $k\ge1$ y $0\le b\le\lfloor k/2\rfloor$ resulta
\[
 N_{k,b}=\frac{2k}{k-b}\binom{k-b}{b}.
\]
En efecto, el coeficiente de $z^kt^b$ en
$(1-z-tz^2)^{-1}$ es $\binom{k-b}{b}$. La multiplicación por el numerador
produce
\[
 3\binom{k-b}{b}-\binom{k-b-1}{b}+\binom{k-b-1}{b-1}
 =2\binom{k-b}{b}+2\binom{k-b-1}{b-1},
\]
que da la expresión anterior. El término constante es $N_{0,0}=3$.
Para $k=4$ hay catorce paseos, distribuidos en $2$, $8$ y $4$ según
contengan cero, uno o dos retornos. El valor publicado no determina
esta composición interna.

\paragraph{Recuperación reversible.}
Se ordenan los sectores por $(n,i,j)$ y, dentro de cada sector, los paseos
lexicográficamente. El número de continuaciones de longitud $m$ desde
$v$ hasta $j$ es $(F^m)_{vj}$. La suma sucesiva de esos cardinales asigna
a cada paseo un índice $a$, con $0\le a<N_k$. Las restas inversas
reconstruyen sucesivamente sus estados. Por inducción en la longitud,
la codificación $C(\gamma)=(k,a)$ y la decodificación $D$ satisfacen
\[
 D C=\mathrm{id},\qquad C D=\mathrm{id}.
\]
El orden lexicográfico serializa la ruta existente; no selecciona una
trayectoria del generador.

Si $A_v$ prolonga por una arista admisible y $T$ elimina el último estado,
sus representaciones $E_v=C A_v D$ y $\widehat T=C T D$ verifican
\[
 \widehat T E_v=\mathrm{id},\qquad
 k(A_v\gamma)-k(\gamma)=1+x_n-v.
\]
La recuperación respeta, por tanto, prolongación y truncación.

Conocido $k$, una representación binaria inyectiva de longitud fija requiere
al menos $\lceil\log_2N_k\rceil$ bits. El índice anterior alcanza esa
longitud. Este mínimo no incluye la transmisión de $k$ ni supone un
código autodelimitado. La capacidad combinatoria satisface
\[
 \lim_{k\to\infty}\frac{\log_2N_k}{k}=\log_2\varphi.
\]
Esta magnitud cuantifica distinción de rutas; su identificación con
una entropía física exigiría el transductor correspondiente.

